% ===================================================================
%  Biologie: Nahrungsnetz als Graph (Knoten = Lebewesen,
%  Pfeil = "wird gefressen von"), angeordnet mit der Bibliothek positioning.
%  Kompilieren: pdflatex vorlage_nahrungsnetz.tex
% ===================================================================
\documentclass[border=2mm]{standalone}
\usepackage[T1]{fontenc}
\usepackage{lmodern}
\usepackage[ngerman]{babel}
\usepackage{tikz}
\usetikzlibrary{arrows.meta, positioning}

\definecolor{akzentrot}{HTML}{DC2626}
\definecolor{akzentgruen}{HTML}{16A34A}
\definecolor{akzentorange}{HTML}{EA580C}

\begin{document}
\begin{tikzpicture}[node distance=1.1cm and 1.6cm, >={Stealth[length=6pt]},
    lw/.style={draw, rounded corners, font=\sffamily\small, minimum width=2cm, minimum height=7mm},
    pflanze/.style={lw, fill=akzentgruen!20}, tier/.style={lw, fill=akzentorange!20},
    raeuber/.style={lw, fill=akzentrot!20}]
  \node[pflanze] (gras) {Gras};
  \node[pflanze, right=of gras] (eiche) {Eiche};
  \node[tier, above=of gras] (maus) {Maus};
  \node[tier, above=of eiche] (raupe) {Raupe};
  \node[tier, right=of raupe] (reh) {Reh};
  \node[tier, above=of raupe] (meise) {Meise};
  \node[raeuber, above=of maus] (fuchs) {Fuchs};
  \node[raeuber, above=of meise] (habicht) {Habicht};
  \foreach \a/\b in {gras/maus, eiche/raupe, eiche/maus, gras/reh, eiche/reh,
                     raupe/meise, maus/fuchs, maus/habicht, meise/habicht}
    \draw[->, thick] (\a) -- (\b);
\end{tikzpicture}
\end{document}
